% Task C17
% Source volume: lower
% Physical scan pages: 149--158
% Printed book pages: 137--146
% Transcribe only source content; omit running headers, footers, and printed page numbers.
\chapter{高维空间中的点集与基本定理}

% BEGIN TRANSCRIPTION

本章的内容是 $\RR$ 中的点集与实数系基本定理在 $\RR^n$ 中的推广，这是研究多元函数的基础。在 §17.1 依照位置关系与密切程度进行点和集合的分类，并讨论其基本性质。§17.2 是 $\RR^n$ 中的基本定理。最后一节是学习要点和两组参考题。

\section{点与点集的定义及其基本性质}

\subsection{点的分类及其性质}

\paragraph{1. 内点、外点、边界点}
先回忆一下 $\RR^n$ 中的距离与邻域的定义。我们知道点
\[
  x=(x_1,x_2,\ldots,x_n)\in\RR^n
\]
的 Euclid 范数（又称模）$\abs{x}$ 定义为
\[
  \abs{x}=\left(\sum_{i=1}^n x_i^2\right)^{1/2}.
\]
由此可引进 $\RR^n$ 中任意两点 $x$ 与 $y$ 的 Euclid 距离为
\[
  d(x,y)=\abs{x-y}
  =\left[\sum_{i=1}^n(x_i-y_i)^2\right]^{1/2}.
\]
与距离有关的最重要的不等式是三角形不等式（参见上册第 6 页）：
\[
  \abs{x-z}\leq \abs{x-y}+\abs{y-z}.
\]
和 $\RR$ 中的邻域定义相仿，可通过距离定义 $\RR^n$ 中的邻域。设点 $a\in\RR^n$，$\delta>0$，称
\[
  O_\delta(a)=\{x\in\RR^n\mid \abs{x-a}<\delta\}
\]
是点 $a$ 的 $\delta$ 邻域，也称其为以点 $a$ 为中心，以 $\delta$ 为半径的 $n$ 维开球。

在 $\RR^n$ 中给定一个集合 $S$，按照点与集合 $S$ 的位置关系可将 $\RR^n$ 中的点分为三类：$S$ 的内点、外点、边界点。具体地说，对于 $\RR^n$ 中的某一点 $x$，若存在它的一个邻域 $O_\delta(x)\subset S$，则称 $x$ 为 $S$ 的内点；若存在 $x$ 的一个邻域 $O_\delta(x)\cap S=\varnothing$，则称 $x$ 为 $S$ 的外点；若在 $x$ 的任一邻域中既有属于 $S$ 的点，又有不属于 $S$ 的点，则称 $x$ 为 $S$ 的边界点。

$S$ 的全体内点组成的集合称为 $S$ 的内部，记为 $\operatorname{int}S$ 或 $S^\circ$。

$S$ 的全体边界点组成的集合称为 $S$ 的边界，记为 $\partial S$。

\paragraph{2. 聚点}
上述分类是按照任一点 $x\in\RR^n$ 的邻域内的点是否属于 $S$ 来进行的。如果按照去心邻域进行分类，则可将 $\RR^n$ 中的点分为 $S$ 的聚点与非聚点两大类。确切地说，对于点 $x\in\RR^n$，如果在 $x$ 的任一去心邻域中总有 $S$ 的点，则称 $x$ 为 $S$ 的聚点。$S$ 的全体聚点组成的集合记为 $S^d$，称为 $S$ 的导集。显然内点一定是聚点，外点一定不是聚点。

如果点 $x\in S$，且存在 $x$ 的一个邻域 $O_\delta(x)\cap S=\{x\}$，则称 $x$ 为 $S$ 的孤立点。孤立点一定不是聚点，而边界点有可能是聚点也有可能是孤立点。

聚点是一个重要概念，它的下述两个等价定义是经常要用到的。

\begin{xdefinition}
\textbf{（1）} 设点 $x\in\RR^n$，如果在它的任何邻域 $O_\delta(x)$ 内总含有 $S$ 中的无穷多个点，则称 $x$ 是 $S$ 的一个聚点。
\end{xdefinition}

\begin{xdefinition}
\textbf{（2）} 设点 $x\in\RR^n$，如果存在由相异点组成的一个点列 $\{x_n\}\subset S$，$x_n\neq x$（$n=1,2,\ldots$），使得 $x_n\to x$，则称 $x$ 为 $S$ 的一个聚点，这里 $x_n\to x$ 的含义是 $d(x_n,x)\to0$。
\end{xdefinition}

\begin{xexample}[17.1.1]
证明：集合 $S$ 的导集的聚点是 $S$ 的聚点，即 $(S^d)^d\subset S^d$。
\end{xexample}
\begin{xproof}
设点 $x\in(S^d)^d$，则 $\exists S$ 的相异聚点 $x_n$（$n=1,2,\ldots$），$x_n\neq x$，且 $x_n\to x$。从而 $\forall\delta>0$，$\exists N$，当 $n>N$ 时，$x_n\in O_\delta(x)$。设 $n_0>N$，由于 $x_{n_0}$ 为 $S$ 的聚点，于是在 $O_{\delta-\abs{x-x_{n_0}}}(x_{n_0})$ 中含有无穷多个 $S$ 中异于 $x_{n_0}$ 的点。显然
\[
  O_{\delta-\abs{x-x_{n_0}}}(x_{n_0})\subset O_\delta(x),
\]
所以 $O_\delta(x)$ 中有无穷多个异于 $x$ 的 $S$ 中的点，由等价定义（1）知 $x$ 为 $S$ 的聚点。
\end{xproof}

\subsection{集合的分类及其性质}

\paragraph{1. 开集、闭集}
如果 $\operatorname{int}S=S$，则称 $S$ 为开集。开集有如下重要性质：
\begin{enumerate}[label=(\arabic*)]
  \item 任意多个开集的并集是开集；
  \item 有限多个开集的交集是开集；
  \item 全空间 $\RR^n$ 和空集 $\varnothing$ 都是开集。
\end{enumerate}

开集的余集定义为闭集。又定义 $S$ 的闭包 $\overline S$ 为
\[
  \overline S=S\cup S^d.
\]
易证 $\overline S$ 为闭集，且 $\overline S=S\cup\partial S$。关于闭集，下列条件等价：
\begin{enumerate}[label=(\arabic*)]
  \item $S$ 是闭集；
  \item $S^d\subset S$（即 $S=\overline S$）；
  \item $\partial S\subset S$（即 $S=\overline S$）。
\end{enumerate}

\begin{xexample}[17.1.2]
设 $S$ 为 $\RR^n$ 中的一个集合，则 $\partial S$ 为闭集。
\end{xexample}
\begin{xproof}
\textbf{1（按定义证）} 假设点 $x\in(\partial S)^c$，即 $\partial S$ 的余集，则 $x$ 只能是 $S$ 的内点或外点。

若点 $x\in\operatorname{int}S$，则 $\exists\delta>0$，使得 $O_\delta(x)\subset S$，由内点定义知 $O_\delta(x)\subset\operatorname{int}S$，从而 $O_\delta(x)\subset(\partial S)^c$；

若 $x$ 是 $S$ 的外点，则 $\exists\delta>0$，使得 $O_\delta(x)\cap S=\varnothing$，因此 $O_\delta(x)\subset S^c$，而 $O_\delta(x)$ 本身是开集，这说明 $O_\delta(x)$ 中的点都不是 $S$ 的边界点，即 $O_\delta(x)\subset(\partial S)^c$。

由定义知 $(\partial S)^c$ 为开集，即 $\partial S$ 为闭集。
\end{xproof}
\begin{xproof}
\textbf{2（证 $(\partial S)^d\subset\partial S$）} 设点 $x\in(\partial S)^d$，由聚点的等价定义（2）知存在相异的点列 $\{x_n\}\subset\partial S$，$x_n\neq x$，$n=1,2,\ldots$，使得 $x_n\to x$，于是 $\forall\delta>0$，$\exists N$，当 $n>N$ 时点 $x_n\in O_\delta(x)$，取 $n_0>N$，由于点 $x_{n_0}\in\partial S$，则由边界点的定义知 $O_{\delta-\abs{x_{n_0}-x}}(x_{n_0})\subset O_\delta(x)$ 中有 $S$ 中的点，也有不在 $S$ 中的点，所以 $x\in\partial S$。
\end{xproof}
\begin{xproof}
\textbf{3（证 $\partial(\partial S)\subset\partial S$）} 设点 $x\in\partial(\partial S)$，则 $\forall\delta>0$，在 $O_{\delta/2}(x)$ 中有 $\partial S$ 的点 $y$。又由边界点定义，在 $O_{\delta/2}(y)$ 中既有属于 $S$ 的点，也有不属于 $S$ 的点。由于 $O_{\delta/2}(y)\subset O_\delta(x)$，因此 $O_\delta(x)$ 中既有属于 $S$ 的点，也有不属于 $S$ 的点，于是 $x\in\partial S$。
\end{xproof}

\paragraph{2. 紧集、凸集}
设 $S$ 是 $\RR^n$ 的一个集合，如果在 $S$ 的任何一个无限开覆盖 $\{O_\alpha\}_{\alpha\in I}$ 中总可以找出有限个开集 $O_1,O_2,\ldots,O_k$，同样可以覆盖 $S$，即
\[
  \bigcup_{i=1}^k O_i\supset S,
\]
则称 $S$ 是 $\RR^n$ 的一个紧集。容易证明紧集一定是有界闭集，而且我们将会看到，在 $\RR^n$ 中紧集与有界闭集的定义是等价的（紧性定理）。

设 $E$ 是 $\RR^n$ 的一个集合，若 $\forall x_1,x_2\in E$，有 $x=tx_1+(1-t)x_2\in E$（$0\leq t\leq1$），则称 $E$ 为凸集。从几何上看，以 $x_1,x_2$ 为端点的直线段位于 $E$ 内。

\begin{xexample}[17.1.3]
紧集的闭子集是紧集。
\end{xexample}
\begin{xproof}
设 $E$ 是一个紧集，$F$ 是 $E$ 的闭子集。设 $\{O_\lambda\}_{\lambda\in I}$ 是 $F$ 的任一开覆盖，由于 $F^c=\RR^n-F$ 是开集，则 $\{O_\lambda\}_{\lambda\in I}$ 与 $F^c$ 一起形成紧集 $E$ 的一个开覆盖。由紧集的定义知在 $\{O_\lambda\}_{\lambda\in I}$ 与 $F^c$ 中存在有限个开集形成 $E$ 的一个有限覆盖，记这有限个开集为 $O_1,O_2,\ldots,O_k$，不妨设 $F^c=O_k$。由于 $F\subset E$，则
\[
  \left(\bigcup_{i=1}^{k-1}O_i\right)\cup F^c\supset E\supset F.
\]
但 $F^c\cap F=\varnothing$，所以
\[
  \bigcup_{i=1}^{k-1}O_i\supset F.
\]
由紧集的定义知 $F$ 为紧集。
\end{xproof}

\paragraph{3. 连通集、区域}
设 $D$ 是 $\RR^n$ 的一个集合，如果当 $D$ 分解为两个不相交的非空子集的并集 $A\cup B$ 时，有 $A^d\cap B\neq\varnothing$ 或者 $A\cap B^d\neq\varnothing$，则称 $D$ 为连通集。当 $D$ 是开集时，我们有：开集 $D$ 是连通集的充分必要条件是 $D$ 不能分解为两个不相交的非空子开集的并（第二组参考题 1）。在 $\RR$ 中，连通集有特别直观的描述：$\RR$ 中集合 $D$ 是连通集的充分必要条件是 $D$ 为区间（第二组参考题 2）。

连通的开集称为区域或开区域。开区域的闭包称为闭区域。

更为直观并易于判断的概念是道路连通集。设 $D$ 是 $\RR^n$ 的一个集合，如果当 $D$ 内任何两点 $p,q$，都可以找到连续曲线 $l\subset D$ 将 $p$ 和 $q$ 联结，则称 $D$ 为道路连通集。这里的连续曲线是指 $l$ 可以表示为参数方程
\[
  x_i=\varphi_i(t),\qquad i=1,2,\ldots,n,
\]
其中诸 $\varphi_i$ 是区间 $[0,1]$ 上的连续函数，并且
\[
  p=(\varphi_1(0),\varphi_2(0),\ldots,\varphi_n(0)),\qquad
  q=(\varphi_1(1),\varphi_2(1),\ldots,\varphi_n(1)).
\]
可以证明道路连通集一定是连通集，但连通集未必是道路连通集（第二组参考题 5）。下面的命题说明了区域的道路连通性。

\begin{xproposition}[17.1.1]
$\RR^n$ 中的区域都是道路连通的。
\end{xproposition}
\begin{xproof}
设 $D$ 是 $\RR^n$ 中的一个非空连通开集。取点 $x\in D$，设 $U(x)$ 为 $D$ 中所有与 $x$ 有 $D$ 中连续曲线相联结的点的集合。容易看到 $U(x)$ 是一个道路连通集。我们证明 $U(x)=D$。设 $y\in U(x)$，并取 $\delta>0$ 使 $O_\delta(y)\subset D$。$\forall z\in O_\delta(y)$ 存在 $O_\delta(y)$ 中的直线段联结 $z$ 到 $y$，从而存在 $D$ 中的连续曲线联结 $z$ 到 $x$，所以 $O_\delta(y)\subset U(x)$。因而 $U(x)$ 是包含 $x$ 的开集。如 $D-U(x)\neq\varnothing$，则 $D-U(x)=\bigcup U(y')$，其中 $y'$ 取遍 $D-U(x)$ 中的点。按前面的证明，每个 $U(y')$ 都是开集，因此 $D-U(x)$ 也是开集。$D$ 有开集分解式
\[
  D=U(x)\cup(D-U(x)),
\]
与 $D$ 是连通开集矛盾。这就证明了 $D-U(x)=\varnothing$，所以 $D=U(x)$ 是道路连通集。
\end{xproof}

\paragraph{4. 距离概念的推广}
点与点的距离概念可推广到点 $x$ 与集合 $S$，集合 $S_1$ 与集合 $S_2$ 之间的距离：
\[
  d(x,S)=\inf_{y\in S}\{\abs{x-y}\};
\]
\[
  d(S_1,S_2)
  =\inf_{\substack{x\in S_1\\y\in S_2}}\{\abs{x-y}\}
  =\inf_{x\in S_1}\{d(x,S_2)\}
  =\inf_{y\in S_2}\{d(y,S_1)\}.
\]
点与集合的距离可以看作是两个集合之间的距离的特殊情况。同时还可以定义一个集合 $S$ 的直径 $d_S$ 为
\[
  d_S=\sup_{x,y\in S}\{\abs{x-y}\}.
\]

关于集合的运算有下列命题：

\begin{xproposition}[17.1.2（De Morgan（德摩根）法则）]
设 $A_\alpha$（$\alpha\in I$）为一族集合，则有
\[
  \left(\bigcup_{\alpha\in I}A_\alpha\right)^c
  =\bigcap_{\alpha\in I}(A_\alpha)^c,
  \qquad
  \left(\bigcap_{\alpha\in I}A_\alpha\right)^c
  =\bigcup_{\alpha\in I}(A_\alpha)^c.
\]
\end{xproposition}

由命题 17.1.2 易证下述结论：任意多个闭集的交集仍是闭集；有限多个闭集的并集仍是闭集。

\subsection{思考题}
\begin{xthought}
\begin{enumerate}[label=\arabic*.]
  \item 按定义证明闭集的如下重要性质：
  \begin{enumerate}[label=(\arabic*)]
    \item 任意多个闭集的交集是闭集；
    \item 有限多个闭集的并集是闭集。
  \end{enumerate}
  \item 证明聚点定义（1）、（2）的等价性。
  \item 在例题 17.1.1 的证明中，我们使用的是聚点等价定义（1）。若使用原始定义，证明是否能通过？若不能，应如何修改？
  \item 无限多个开集的交是否一定是开集？
\end{enumerate}
\end{xthought}

\subsection{练习题}
\begin{xexercise}
\begin{enumerate}[label=\arabic*.]
  \item 证明：$\overline S=S\cup\partial S$。
  \item 证明：$\partial S=\overline S-\operatorname{int}S$。
  \item 若 $A\cap B=\varnothing$，则 $\overline A\cap(\operatorname{int}B)=\varnothing$。
  \item 证明：$S=S^d\Longleftrightarrow S$ 闭，且 $S$ 无孤立点。
  \item $S$ 为 $\RR^n$ 中的点集，证明：$\overline S=\{x\in\RR^n\mid d(x,S)=0\}$。
  \item 若 $S$ 为凸集，则 $\overline S$ 也是凸集。
  \item 对于集合 $S$ 与任一组集合 $A_\alpha$，$\alpha\in I$，恒有分配律：
  \[
    S\cap\left(\bigcup_{\alpha\in I}A_\alpha\right)
    =\bigcup_{\alpha\in I}(S\cap A_\alpha).
  \]
\end{enumerate}
\end{xexercise}

\section{$\RR^n$ 中的几个基本定理}

\subsection{综述}

$\RR$ 中的六个基本定理（见上册第三章）能推广到 $\RR^n$（$n>1$）上的是四个定理，它们是：
\begin{enumerate}[label=(\arabic*)]
  \item 闭矩形套定理；
  \item 凝聚定理：$\RR^n$ 中的有界点列一定有收敛子列（或聚点定理：有界无限点集一定有聚点）；
  \item Cauchy 收敛准则：收敛点列 $\Longleftrightarrow$ 基本点列；
  \item 紧性定理：$\RR^n$ 中的点集 $S$ 是紧集的充分必要条件是 $S$ 为有界闭集（覆盖定理）。
\end{enumerate}
其他两个定理（确界存在定理，单调有界定理）之所以不能推广到高维空间，是因为它们与一维直线上的点的顺序有关。

紧性定理的叙述与一维的覆盖定理不同，这可以从两方面进行解释：
\begin{enumerate}[label=(\arabic*)]
  \item 如果在一维的情况下我们也定义闭集与紧集，则覆盖定理就叙述为：有界闭区间是紧集（参见上册 80--81 页）；
  \item 一维的覆盖定理不能以充分必要条件的形式叙述，因为那时没有定义闭集，而一维紧集是有界闭集但不一定是有界闭区间。
\end{enumerate}
下面我们利用 De Morgan 法则给出紧性定理的另一种等价的表达形式。

\begin{xdefinition}
设集合 $S\subset\RR^n$，称 $\RR^n$ 中的子集族 $\{F_\lambda\}_{\lambda\in I}$ 关于 $S$ 具有有限交性质，若对于 $I$ 的任何有限子集 $J$ 均有
\[
  S\cap\left(\bigcap_{\lambda\in J}F_\lambda\right)\neq\varnothing.
\]
\end{xdefinition}

\begin{xproposition}[17.2.1]
$\RR^n$ 中的集合 $S$ 是紧集的充分必要条件是任何关于 $S$ 具有有限交性质的闭集族 $\{F_\lambda\}_{\lambda\in I}$ 与 $S$ 必有非空交，即
\[
  S\cap\left(\bigcap_{\lambda\in I}F_\lambda\right)\neq\varnothing.
\]
\end{xproposition}
\begin{xproof}
先证充分性。设任一关于 $S$ 具有有限交性质的闭集族与 $S$ 有非空交。任取 $S$ 的一个开覆盖 $\{O_\lambda\}_{\lambda\in I}$，则由 De Morgan 法则，从 $\bigcup_{\lambda\in I}O_\lambda\supset S$ 得
\[
  \bigcap_{\lambda\in I}O_\lambda^c\subset S^c,
\]
即
\[
  S\cap\left(\bigcap_{\lambda\in I}O_\lambda^c\right)=\varnothing.
\]
由条件知，$\bigcap_{\lambda\in I}O_\lambda^c$ 关于 $S$ 无有限交性质，即存在有限个 $O_1^c,O_2^c,\ldots,O_k^c$，使得
\[
  S\cap\left(\bigcap_{i=1}^kO_i^c\right)=\varnothing,
\]
从而
\[
  \bigcap_{i=1}^kO_i^c\subset S^c.
\]
于是
\[
  \bigcup_{i=1}^kO_i\supset S.
\]
这样的 $O_i$（$i=1,2,\ldots,k$）就是 $S$ 的一个有限开覆盖，所以 $S$ 为紧集。

再证必要性。设 $S$ 为紧集，$\{F_\lambda\}_{\lambda\in I}$ 是任一关于 $S$ 具有有限交性质的闭集族。假设
\[
  S\cap\left(\bigcap_{\lambda\in I}F_\lambda\right)=\varnothing,
\]
即可由 De Morgan 法则推出矛盾，从略。
\end{xproof}

\subsection{例题}

\begin{xexample}[17.2.1（闭集套定理）]
设 $\{D_k\}$ 是一列非空闭集，它满足：
\begin{enumerate}[label=(\arabic*)]
  \item $D_{k+1}\subset D_k$，$k=1,2,\ldots$；
  \item $D_k$ 的直径 $\delta_k=\displaystyle\sup_{x,y\in D_k}\{\abs{x-y}\}\to0$（$k\to\infty$），
\end{enumerate}
则这列闭集 $D_k$（$k=1,2,\ldots$）存在唯一的公共点。
\end{xexample}
\begin{xproof}
\textbf{1（用凝聚定理）} 在每个 $D_k$ 中任取一点 $x_k$，则 $\{x_k,k=1,2,\ldots\}$ 为一有界无穷点列，由凝聚定理存在点 $x$ 及 $\{x_k\}$ 的子列 $\{x_{k_i},i=1,2,\ldots\}$ 使得
\[
  x_{k_i}\to x\qquad(i\to\infty).
\]
下面证 $x$ 为 $D_k$（$k=1,2,\ldots$）的公共点。事实上，$\forall k_0\in\NN_+$，当 $k_i\geq k_0$ 时
\[
  x_{k_i}\in D_{k_i}\subset D_{k_0}.
\]
令 $i\to\infty$，由于 $D_{k_0}$ 是闭集，则 $x\in D_{k_0}$。

下证唯一性。用反证法，若存在两个公共点 $x,x^*$，记 $d=\abs{x-x^*}$，则 $d>0$，由于 $\delta_k\to0$，于是 $\exists K$，当 $k>K$ 时，$\delta_k<d$，此与 $x,x^*\in D_k$ 矛盾。
\end{xproof}
\begin{xproof}
\textbf{2（用 Cauchy 收敛准则）} 在每个 $D_k$ 中取一点 $x_k$，则
\[
  \abs{x_k-x_l}\leq\max\{\delta_k,\delta_l\}.
\]
由 Cauchy 收敛准则知存在点 $x$，使得 $x_k\to x$。又 $\forall k_0\in\NN_+$，当 $k>k_0$ 时
\[
  x_k\in D_k\subset D_{k_0}.
\]
令 $k\to\infty$，则 $x\in D_{k_0}$，即 $x$ 为 $D_k$（$k\geq1$）的公共点。唯一性的证明同证 1。
\end{xproof}

\begin{xexample}[17.2.2]
设 $S$ 为 $\RR^n$ 中的集合，若 $S$ 既开且闭，则 $S=\RR^n$ 或 $S=\varnothing$。
\end{xexample}
\begin{xproof}
\textbf{1} 因 $S$ 是闭集，故 $\RR^n-S$ 是开集。于是 $\RR^n=S\cup(\RR^n-S)$ 是两个不相交开集的并。由 $\RR^n$ 的连通性可知 $S$ 与 $\RR^n-S$ 中至少有一个为空集，故 $S=\RR^n$ 或 $S=\varnothing$。
\end{xproof}
\begin{xproof}
\textbf{2（不用连通性概念的证明）} 首先证明 $\partial S=\varnothing$。因为 $S$ 开，$S$ 内的点都是内点，所以 $S$ 内无 $S$ 的边界点。同理 $\RR^n-S$ 内也无 $S$ 的边界点，因此 $\partial S=\varnothing$。如果 $S$ 与 $\RR^n-S$ 均非空，则存在点 $x\in S$，点 $y\in\RR^n-S$。设 $L$ 是联结 $x$ 与 $y$ 的直线段，则 $L$ 是有界闭集。设 $z$ 是 $L$ 的中点，则 $z\in S$ 或 $z\in\RR^n-S$。因而 $L$ 有子直线段 $L_1$ 分别以 $S$ 与 $\RR^n-S$ 中的点为其端点，依此可构造由 $L$ 的子直线段组成的有界非空闭集套
\[
  L\supset L_1\supset\cdots\supset L_i\supset\cdots,
\]
其集合半径趋于零，且两端点分别为 $S$ 与 $\RR^n-S$ 中的点。由闭集套定理，存在唯一的点 $p$ 属于所有的直线段。由边界点的定义可见 $p\in\partial S$，与 $\partial S=\varnothing$ 矛盾。
\end{xproof}

\begin{xexample}[17.2.3]
按 $(1)\Rightarrow(2)\Rightarrow(3)\Rightarrow(1)$ 的次序证明下列三个命题的等价性。
\begin{enumerate}[label=(\arabic*)]
  \item $S$ 是紧集；
  \item $S$ 的任一无限子集必有聚点在 $S$ 中；
  \item $S$ 是有界闭集。
\end{enumerate}
\end{xexample}
\begin{xproof}
$(1)\Rightarrow(2)$。用反证法：设 $F$ 是 $S$ 的无限子集，$\forall x\in S$，它都不是 $F$ 的聚点（这其中有两种可能，一是 $F$ 没有聚点，一是 $F$ 有聚点但不在 $S$ 中）。由聚点的定义 $\exists\delta_x>0$，使得在 $O_{\delta_x}(x)$ 中没有异于 $x$ 的 $F$ 的点。由于
\[
  \bigcup_{x\in S}O_{\delta_x}(x)\supset S\supset F
\]
以及 $S$ 是紧集，从而存在有限个 $O_{\delta_1}(x_1),O_{\delta_2}(x_2),\ldots,O_{\delta_k}(x_k)$ 满足
\[
  \bigcup_{i=1}^kO_{\delta_i}(x_i)\supset S\supset F.
\]
由此可看出 $F$ 是有限集，与 $F$ 是无限集矛盾。

$(2)\Rightarrow(3)$。由已知条件知 $S^d\subset S$，从而 $S$ 闭。下面用反证法证明 $S$ 有界，若不然在 $S$ 中有子列 $\{x_k\}$ 满足 $\abs{x_k}\to+\infty$（$k\to\infty$），可见 $\{x_k\}$ 没有聚点，与已知条件矛盾。

$(3)\Rightarrow(1)$。用闭矩形套定理。其详细证明在很多教科书中都有，从略。
\end{xproof}

\begin{xexample}[17.2.4]
设 $S_1,S_2$ 都是 $\RR^n$ 中的有界闭集，$S_1\cap S_2=\varnothing$，证明：存在两个开集 $O_1$ 和 $O_2$，使得 $S_i\subset O_i$，$i=1,2$，且 $O_1\cap O_2=\varnothing$。
\end{xexample}

\paragraph{分析}
若 $S_1$ 与 $S_2$ 只是两个单点集 $\{x_1\}$ 与 $\{x_2\}$，则 $d=d(x_1,x_2)>0$。令
\[
  O_i=\{x\in\RR^n\mid d(x,x_i)<d/3\},\qquad i=1,2,
\]
则 $O_i$ 为开集，$O_i\supset S_i$，$i=1,2$，且 $O_1\cap O_2=\varnothing$。由此启发我们对一般的有界闭集把证明分成两部分。

\begin{xproof}
第一步：证明 $d=d(S_1,S_2)>0$。用反证法。若 $d(S_1,S_2)=0$，则由 $d(S_1,S_2)$ 的定义，$\exists x_n\in S_1$，$y_n\in S_2$，使得
\begin{equation}
  \lim_{n\to\infty}\abs{x_n-y_n}=0.
  \tag{17.1}
  \label{eq:17.1}
\end{equation}
由 $S_1,S_2$ 有界及凝聚定理知，$\{x_n\},\{y_n\}$ 都有收敛子列，不妨设 $x_n\to x$，$y_n\to y$。由 $S_1,S_2$ 闭知 $x\in S_1$，$y\in S_2$。在 \eqref{eq:17.1} 中令 $n\to\infty$ 得 $x=y$，此与 $S_1\cap S_2=\varnothing$ 矛盾。

第二步：直接定义
\[
  O_i=\{x\in\RR^n\mid d(x,S_i)<d/3\},\qquad i=1,2,
\]
则 $O_i$ 为开集，$O_i\supset S_i$，$i=1,2$，且 $O_1\cap O_2=\varnothing$。

还有一种更有启发性（可用于下面第一组参考题的第 5 题）的构造 $O_1,O_2$ 的办法是定义
\[
  O_1=\bigcup_{x\in S_1}O_{d_x/3}(x),\qquad
  O_2=\bigcup_{y\in S_2}O_{d_y/3}(y),
\]
其中 $d_x$ 是 $x\in S_1$ 到 $S_2$ 的距离，$d_y$ 是 $y\in S_2$ 到 $S_1$ 的距离。
\end{xproof}

\subsection{练习题}
\begin{xexercise}
\begin{enumerate}[label=\arabic*.]
  \item 设 $A,B$ 是 $\RR^n$ 中的两个不相交的闭集，其中一个有界，证明：$d(A,B)>0$。如果 $A,B$ 均是无界闭集，是否仍有 $d(A,B)>0$？
  \item 设 $S_1,S_2$ 为 $\RR^n$ 中不相交的闭集，其中一个有界。证明：存在开集 $O_1,O_2$ 满足 $S_i\subset O_i$，$i=1,2$，且 $O_1\cap O_2=\varnothing$。
  \item 由闭矩形套定理证明凝聚定理。
  \item 由凝聚定理证明 Cauchy 收敛定理。
  \item 由紧性定理证明聚点定理。
\end{enumerate}
\end{xexercise}

\section{对于教学的建议}

\subsection{学习要点}

1. 本章有很多集合论的知识，对于初学者来说是比较抽象的，但它们却是进一步学习数学的基本语言之一。在教学中最基本的要求是理解各类集合的概念。根据我们的经验，以聚点作为一个切入点来展开一些练习很有效果。此外，举例是理解概念的最有效的手段，通常的教科书中有足够的例子供习题课选用。

2. 在上册第三章中的实数系基本定理中，除去单调有界定理和确界存在定理需要用到实数的有序性外，其余四个基本定理都得到了推广，其中闭区间套定理被推广成比较方便的闭集套定理的形式。

3. 在以后的学习中，用得最多的集合概念是区域。我们给出了它的比较正规的描述。根据命题 17.1.1，也可以用道路连通开集来定义 $\RR^n$ 中的区域，这样可以在现阶段避开连通与道路连通这两个易混淆的概念。

4. 多元基本定理延续了第三章的内容，相对而言，学生理解不是很困难。因而有条件时可适当布置一些应用题，如本章的第一组参考题 2。

5. 对习题课的建议 教师应该根据实际教学课时来选取材料，确定教案。一般来说，本章的教学时段只有一次多一点的习题课。因此，习题课的安排首先要保证学生能理解概念，能正确叙述多元基本定理的内容；其次是一些初步的集合证明题和基本定理的应用题（可集中于某一个定理，如闭集套定理的应用）。至于连通性的讨论和紧性的等价描述完全是补充内容。

学生们在证明某个集合 $S$ 具有某种性质时，有时采用如下错误证法：先对 $S$ 是开集时证明结论，然后对 $S$ 是闭集时证明结论。由以上两步得到结论成立。要告诉他们尽管开集的余集是闭集，闭集的余集是开集，但这并不意味着集合只有开集和闭集两大类。

在证明像 17.2.3 小节中的练习题 1 时，学生们往往不会下手，或者是不会做，或者是不得要领地写一大堆。要引导他们从最简单的情况开始思考。如果 $A$ 为单点集 $\{x\}$ 时，用反证法，设 $d(x,B)=0$，则存在 $y_n\in B$，使得 $\abs{x-y_n}\to0$（$n\to\infty$），因而 $\{y_n\}$ 是有界列，有收敛子列，然后如例题 17.2.4 第一步那样推出矛盾。当 $A$ 为一般有界闭集时，讨论是类似的（此时如何证明 $\{y_n\}$ 也是有界列呢？）。

\subsection{参考题}

\subsubsection*{第一组参考题}
\begin{xreferenceproblem}
\begin{enumerate}[label=\arabic*.]
  \item 证明：$\RR^n$ 中每个闭集可表为可列个开集的交，每个开集可表为可列个闭集的并。
  \item 用闭集套定理证明三角形三边上的中线交于一点。
  \item 设函数 $f(x)$ 在 $(-\infty,+\infty)$ 上连续，证明：$E=\{(x,f(x))\mid x\in(-\infty,+\infty)\}$ 是平面上的闭集（称为 $f$ 的图像）。
  \item 证明：$\operatorname{int}S$ 是开集，而且是包含于 $S$ 的最大开集。
  \item （$\RR^n$ 的正规性）设 $S_1,S_2$ 为 $\RR^n$ 中不相交的闭集（不一定有界）。证明：存在开集 $O_1,O_2$ 满足 $S_i\subset O_i$，$i=1,2$，且 $O_1\cap O_2=\varnothing$。
  \item
  \begin{enumerate}[label=(\arabic*)]
    \item 设 $S$ 是 $\RR^n$ 中的有界集合。证明：$\forall\delta>0$，$\exists S$ 中有限个点 $p_1,p_2,\ldots,p_k$，使得
    \[
      \bigcup_{i=1}^kO_\delta(p_i)\supset S;
    \]
    \item 证明覆盖定理的 Lebesgue 形式：若 $\{G_\alpha\}$ 是有界闭集 $F$ 的开覆盖，则 $\exists\delta>0$，使得 $\forall p\in F$，存在 $\{G_\alpha\}$ 中的一个开集 $G$，满足 $O_\delta(p)\subset G$；
    \item 证明覆盖定理的 Lebesgue 形式与通常的表达形式（若 $\{G_\alpha\}$ 是有界闭集 $F$ 的开覆盖，则存在 $\{G_\alpha\}$ 中的有限个开集 $G_1,G_2,\ldots,G_m$，使得 $\bigcup_{i=1}^mG_i\supset F$）等价。
  \end{enumerate}
\end{enumerate}
\end{xreferenceproblem}

\subsubsection*{第二组参考题}
\begin{xreferenceproblem}
\begin{enumerate}[label=\arabic*.]
  \item 证明：$\RR^n$ 中开集 $D$ 是连通集的充分必要条件是 $D$ 不能分解为两个不相交的非空子开集的并。
  \item 证明：$\RR$ 中集合 $D$ 是连通集的充分必要条件是 $D$ 为区间。
  \item 证明：有公共点的连通集的并集是连通集。
  \item 证明：$A=\{(x,y)\mid x,y\text{ 至少有一个是有理数}\}$ 是 $\RR^2$ 中的道路连通集。
  \item 证明：道路连通集一定是连通集，但连通集未必是道路连通集。（考察例子
  \[
    E=\{(x,y)\mid y=\sin\tfrac1x,\ 0<x\leq\tfrac2\pi\},
  \]
  证明 $\overline E$ 是连通而非道路连通的。）
  \item 设 $A,B$ 为 $\RR^n$ 的非空子集，定义 $A+B$ 为 $\RR^n$ 中一切形如 $x+y$ 的点组成的集合，其中 $x\in A$，$y\in B$。
  \begin{enumerate}[label=(\arabic*)]
    \item 如果 $K$ 为 $\RR^n$ 的紧子集，$C$ 为 $\RR^n$ 的闭子集，证明：$K+C$ 为 $\RR^n$ 的闭子集；
    \item 如果 $K,C$ 均为 $\RR^n$ 的闭子集，$K+C$ 未必为 $\RR^n$ 的闭子集。考虑 $K$ 为整数集合，$C$ 为一切形如 $\alpha m$ 的数组成的集合，其中 $m\in K$，$\alpha$ 是某个确定的无理数。证明：$\overline{K+C}=\RR$，但 $K+C\neq\RR$。
  \end{enumerate}
\end{enumerate}
\end{xreferenceproblem}
